\newpage

\markright{\textit{Supplementary Note 2}}

\noindent{\fontsize{12}{14.4}\selectfont\textbf{Supplementary Note 2}}

\vspace{0.5em}

\noindent \textbf{Covariate association screening and covariate-adjusted models:} Associations between covariates and individual chemical features were evaluated using the same exposome-wide approach applied to tumor type, so that the rate of covariate association could be benchmarked directly against the rate observed for tumor type itself. Age, sex, and collection year were each screened against every annotated feature in both analytical modes, with collection year evaluated both as a continuous measure and as the categorical intervals used for cohort description. Because the quantitative test is the \textit{F} test from a linear model, the same specification yields a one-way analysis of variance for categorical covariates and a test of the regression slope for continuous covariates, so no additional test-selection decision was introduced; for a single continuous predictor, $\eta^{2}$ is equivalent to $r^{2}$. For qualitative features, Fisher's exact test has no continuous analogue, and a logistic-regression likelihood-ratio test was therefore used for continuous covariates. Effect size (median $\eta^{2}$) was calculated for quantitative features only, because the qualitative tests yield different metrics (Cramér's \textit{V} for categorical covariates and McFadden's pseudo-$R^{2}$ for continuous covariates) that are not directly comparable to one another. Because the inference concerns the proportion of features associated with a covariate relative to the approximately 5\% expected under the null, rather than the identity of individual features, no per-feature multiplicity correction was applied and no individual covariate-feature associations are reported.

\vspace{0.6em}

\noindent For covariate adjustment, each chemical showing a tumor-type difference was refit under three nested models. A baseline model (“Model 1”) tested the tumor-type term alone, using a linear-model \textit{F} test for quantitative features and a logistic-regression likelihood-ratio test for qualitative features. This deviates from the exposome-wide analysis in the qualitative mode only: because Fisher's exact test cannot accommodate covariates, the likelihood-ratio test was used throughout so that unadjusted and adjusted models remain directly comparable. Quantitative results are unaffected, as the one-way analysis of variance and the linear-model \textit{F} test are the same test. Qualitative \textit{P} values therefore differ marginally from the corresponding exact tests, although the same chemicals were significant under both. A second model (“Model 2”) added adjustment for collection year, and a third (“Model 3”) adjusted for collection year, age, and sex. Effect size was partial $\eta^{2}$ for quantitative features and McFadden's pseudo-$R^{2}$ for qualitative features; because these metrics are on different scales, they should be compared only within a mode and never across modes. Adjusted estimates for features detected in ten or fewer samples should be interpreted with caution, given the limited number of events available per model parameter. A parallel sensitivity analysis substituting the categorical collection intervals for continuous collection year was also fit.
